\section{Sum of Divisors}
\label{sec:n-sum-div}
\source{CSES 1082 (Notion: Number Theory)}{Medium}

\problemstatement{Compute $\sum_{i=1}^{n}\sigma(i)\bmod(10^9+7)$ for
$n\le10^{12}$, where $\sigma(i)$ is the sum of divisors of $i$.}

\begin{patternbox}
Swap the order of summation (count how often each divisor appears), then
group equal values of $\lfloor n/d\rfloor$ --- \textbf{floor blocks}.
\end{patternbox}

\subsection*{Approach and intuition}

Divisor $d$ divides exactly $\lfloor n/d\rfloor$ numbers in $[1,n]$, so
\[
  \sum_{i\le n}\sigma(i)=\sum_{d=1}^{n}d\left\lfloor\frac nd\right\rfloor .
\]
$\lfloor n/d\rfloor$ has at most $2\sqrt n=2\cdot10^6$ distinct values. For a
block $[\ell,r]$ where it equals $q$ (with $r=\lfloor n/q\rfloor$), add
$q\cdot(\ell+\dots+r)$, computing the arithmetic series modulo $p$ with the
inverse of $2$.

\subsection*{Algorithm}
\begin{algorithmic}[1]
\State $\ell\gets1$
\While{$\ell\le n$}
  \State $q\gets\lfloor n/\ell\rfloor$; $r\gets\lfloor n/q\rfloor$
  \State $ans\gets ans+q\cdot\frac{(r-\ell+1)(\ell+r)}{2}\pmod p$; $\ell\gets r+1$
\EndWhile
\end{algorithmic}

\dryrun{$n=5$: blocks $[1,1]$ ($q=5$): $5$; $[2,2]$ ($q=2$): $4$; $[3,5]$
($q=1$): $12$. Total $21=1+3+4+7+6$. \checkmark}

\subsection*{C++17 implementation}
\cppfile{number-theory/14_sum_of_divisors.cpp}

\begin{complexitybox}
\textbf{Time:} $O(\sqrt n)$. \textbf{Space:} $O(1)$.
\end{complexitybox}

\begin{pitfallbox}
\begin{itemize}
  \item $(r-\ell+1)(\ell+r)$ reaches $10^{24}$: reduce each factor modulo
        $p$ first, then multiply by $2^{-1}=500000004$.
  \item Reduce $q$ modulo $p$ too before multiplying.
\end{itemize}
\end{pitfallbox}
